\documentclass[numberedappendix]{preprint}

\usepackage{amsmath}
\usepackage{amssymb}
\usepackage{graphicx}
\usepackage{natbib}
\usepackage{hyperref}
\hypersetup{colorlinks=true,linkcolor=blue,citecolor=blue,filecolor=blue,urlcolor=blue}

\newcommand{\uvece}{\hat{\mathbf e}}
\newcommand{\uvecg}{\hat{\mathbf g}}
\newcommand{\uvecr}{\hat{\mathbf r}}
\newcommand{\uvecx}{\hat{\mathbf x}}
\newcommand{\uvecy}{\hat{\mathbf y}}
\newcommand{\uvecz}{\hat{\mathbf z}}
\newcommand{\vecr}{{\mathbf r}}

\begin{document}

\title{AI sans Slop:\\ Basis functions for stellar interferometry\\ developed with AI assistance}
\shorttitle{AI sans Slop}

\author{Perpetua V. Sallow$^{1}$}
\author{Bartholomew K. Nimbus$^{2}$}
\author{Xiomara L. Petrichor$^{3}$}

\affiliation{$^1$Institute of Advanced Procrastination, University of the Empty Quarter, Muscat, Oman}
\affiliation{$^2$Department of Celestial Bureaucracy, University of the Long White Cloud, Wellington, New Zealand}
\affiliation{$^3$Centro de F\'isica del Vac\'io, Universidad del C\'ondor Imaginario, Santiago, Chile}

\shortauthors{Sallow, Nimbus \& Petrichor}

\begin{abstract}
Spin and tide
\end{abstract}

\keywords{methods: miscellaneous --- sociology of astronomy --- artificial intelligence}

\section{Introduction}

Interferometric observations of gravity darkening in stars
\citep[e.g.,][]{2018MNRAS.480.1263H,2018A&A...619A.167D} and the
recent progress of intensity interferometry
\citep{2025ApJ...995..191A} have created interest in basis functions
for stellar interferometry.

\cite{2026PASP..138e4504D} developed such a basis set for a rotating
star.  In this work we will re-derive their basis set with a
refinement that makes it closed-form, and extend it both
perturbatively and numerically to the case of a triaxial star, which
is flattened by spin and stretched by tide in along different axes.
This is done using symbolic computing and AI assistance.

A recent essay by \citet{tao2026} identifies five stages of a new
result in mathematics.
\begin{itemize}
\item \textit{Generation} of new theorem proofs.  Here AI has brought
  revolutionary changes.
\item \textit{Verification} of the proofs.  Symbolic computing had
  made it possible to keep up with the speed of generation.
\item \textit{Exposition} for human readers.  AI has not (yet) helped
  here, and in fact has made things worse --- the problem of AI slop.
\item \textit{Publication} in the research literature.
\item \textit{Canonicalization}, or absorption into textbooks.
\end{itemize}
Here would like to study how the first four stages translate into
astrophysics.

\section{Stellar interferometry}

\begin{equation}
  V(u,v) \propto \int I(x,y)\, e^{-2\pi i (u x + v y)}\; dx\, dy
  \label{eq:visibility-definition}
\end{equation}

\begin{equation}
  V(u,v) \propto \int_{\uvecg\cdot\uvecz > 0}
  I(\vecr\in S)\, 
  e^{-2\pi i (u x + v y)}\; (\uvecg\cdot\uvecz)\, dS
  \label{eq:deformation-visibility-integral}
\end{equation}

\section{Spherical harmonics}

\begin{equation}
  Y_{lm}(\theta,\phi) = N_{lm}\, P_{lm}(\cos\theta)\, e^{im\phi}
\end{equation}
where
\begin{equation}
\begin{aligned}
  P_{lm}(x) &= (1-x^{2})^{m/2}\,\frac{d^{m}}{dx^{m}}P_l(x) \\
  P_l(x)    &= \frac{1}{2^{l}l!}\,\frac{d^{l}}{dx^{l}}\,(x^{2}-1)^{l}
\end{aligned}
\label{eq:associated-legendre}
\end{equation}
and
\begin{equation}
  N_{lm} = \sqrt{(2l+1)\,\frac{(l-m)!}{(l+m)!}} ,
  \label{eq:ylm}
\end{equation}

\begin{equation}
  Y_{l,-m} =   Y_{l,m}^*
\end{equation}

\begin{equation}
\begin{matrix}
    Y_{00} = 1 \hfill
  & Y_{10} = \sqrt{3}\,\cos\theta \hfill
  & Y_{11} = \sqrt{\tfrac32}\,\sin\theta\,e^{i\phi} \hfill \\[4pt]
    Y_{20} = \tfrac{\sqrt{5}}{2}(3\cos^{2}\theta-1) \qquad \hfill
  & Y_{21} = \sqrt{\tfrac{15}{2}}\,\sin\theta\cos\theta\,e^{i\phi}
    \qquad \hfill
  & Y_{22} = \sqrt{\tfrac{15}{8}}\,\sin^{2}\theta\,e^{2i\phi} \hfill
\end{matrix}
\label{eq:ylm-table}
\end{equation}

\def\apj{ApJ}
\def\pasp{PASP}

\bibliographystyle{aasjournal}
\bibliography{bibliography}

\appendix
\section{A note on the void}

A more detailed account of the void will be presented elsewhere.


\end{document}
